% ECONOMICS_PROBLEM: {"id": "F18-493", "title": "Carbon border adjustments and leakage", "jel_code": "F18", "source_id": "OP-0087"}
% !TeX program = xelatex
\documentclass[11pt,a4paper]{article}
\usepackage[margin=23mm]{geometry}
\usepackage{fontspec}
\setmainfont{Latin Modern Roman}
\usepackage{amsmath,amssymb,mathtools}
\usepackage{xcolor,enumitem,booktabs,longtable,array,xurl,hyperref}
\definecolor{muted}{HTML}{53636C}
\hypersetup{colorlinks=true,urlcolor=blue,linkcolor=blue}
\setlength{\parindent}{0pt}
\setlength{\parskip}{5pt}
\newcommand{\R}{\mathbb R}
\newcommand{\E}{\mathbb E}
\newcommand{\N}{\mathbb N}
\newcommand{\ind}{\mathbf 1}
\newcommand{\Var}{\operatorname{Var}}
\newcommand{\Cov}{\operatorname{Cov}}
\newcommand{\supp}{\operatorname{supp}}
\DeclareMathOperator*{\argmin}{arg\,min}
\DeclareMathOperator*{\argmax}{arg\,max}
\newcommand{\entrymeta}[3]{{\sffamily\small\color{muted}\textbf{\texttt{#1}}\quad|\quad #2\par #3\par}}
\newcommand{\Problemref}[1]{\texttt{#1}}
\newcommand{\JELref}[2]{\texttt{#2} (JEL \texttt{#1})}
\begin{document}
\section*{Carbon border adjustments and leakage}
\textbf{Problem ID:} \texttt{F18-493}\quad \textbf{JEL:} \texttt{F18}\par
% BEGIN ECONOMICS STATEMENT
\label{problem:OP-0087}
\label{atlas:prob:T7}

\noindent\textbf{Atlas ID:} T7; \textbf{original number:} 87; \textbf{field:} International trade. \textbf{Classification:} Editorial research formulation (theoretical/empirical); not a certified unsolved theorem.

\subsection*{Definitions and assumptions}
Consider countries $j=1,\ldots,J$ and goods $g=1,\ldots,G$. Firms choose production location, inputs and abatement within specified technologies, generating physical greenhouse-gas emissions $E_j\ge0$ in common carbon-dioxide-equivalent units. Policy $b$ consists of domestic carbon prices, border assessments, exemptions, verified intensity rules and recycling of receipts. The equilibrium selection includes consumer budget optimization, firm profit maximization, market clearing and consistency between verified reports and an explicit measurement or enforcement model. Let $E_j(b)$ denote equilibrium territorial emissions; assume finite expected emissions and utility.

\subsection*{Formal research question}
For a domestic-pricing baseline $b^0$ and border-adjustment policy $b^1$, define global mitigation and, when its denominator is strictly positive, a leakage ratio:
\[
 M=\sum_j\mathbb E[E_j(b^0)-E_j(b^1)],\quad
 \lambda=\frac{\sum_{j\notin D}\mathbb E[E_j(b^1)-E_j(b^0)]}
                 {\sum_{j\in D}\mathbb E[E_j(b^0)-E_j(b^1)]},
\]
where $D$ is the implementing coalition. Characterize implementable border rules that achieve $M>0$ and acceptable household equivalent-variation losses under bounded reporting error and specified administrative-resource costs. Identify substitution, relocation and reporting responses using firm-product data, distinguishing physical emissions from reported embedded emissions.

\subsection*{Interpretation and scope}
The ratio is undefined when domestic reductions vanish and can be negative if foreign emissions also fall. It is not a welfare index. Any optimization over mitigation and welfare must state social weights or the monetary valuation of climate damages; do not add tonnes and utility without conversion. Domestic production, consumption-based carbon footprints and global territorial emissions use different accounting boundaries and should be reconciled. A tariff-equivalent carbon charge can redirect trade without lowering global emissions. Policy-event identification must address concurrent ETS changes, free allocation and reporting adjustments. The implementing coalition and sector coverage must be fixed before calculating the ratio. Embodied-emissions estimates should report a measurement-error distribution or a justified deterministic error bound.

\subsection*{Source audit and status}
[S1] is verified leakage-measurement background. The atlas’s undated ``Böhringer et al.'' and ``carbon-leakage literature'' are incomplete citations; [S2] is a verified relevant candidate, not a uniquely resolved source. [S3] verifies CBAM context and quarterly certificate-price averaging in 2026. These institutional rules do not establish policy effectiveness or certify a presently unresolved theorem.

\subsection*{References}

\begin{enumerate}[label={[S\arabic*]},leftmargin=*]

\item Meredith Fowlie and Mar Reguant (2018). \emph{Challenges in the Measurement of Leakage Risk}. AEA Papers and Proceedings 108, 124–129. \url{https://doi.org/10.1257/pandp.20181087}.
\textit{Locator and role:} Publisher or author abstract / bibliographic record; Definition and measurement challenges of emissions leakage. Provides background or a benchmark result; does not establish that the editorial question remains unresolved in 2026.

\item Christoph Böhringer, Edward J. Balistreri, and Thomas F. Rutherford (2012). \emph{The Role of Border Carbon Adjustment in Unilateral Climate Policy: Overview of an Energy Modeling Forum Study (EMF 29)}. Energy Economics 34, supplement 2, S97–S110. \url{https://doi.org/10.1016/j.eneco.2012.10.003}.
\textit{Locator and role:} Publisher or author abstract / bibliographic record; Verified relevant model-comparison background, supplied editorially. The original author-only citation cannot be uniquely resolved; this is a clearly identified relevant candidate, not a confirmed attribution of the atlas citation.

\item European Commission (2026). \emph{CBAM Definitive Regime}. Directorate-General for Taxation and Customs Union, institutional guidance. \url{https://taxation-customs.ec.europa.eu/carbon-border-adjustment-mechanism/cbam-definitive-regime_en}.
\textit{Locator and role:} Publisher or author abstract / bibliographic record; Certificate prices linked to ETS auctions; quarterly average in 2026, weekly from 2027. Verifies policy context, not demonstrated leakage reduction or the unresolved status of the research question; guidance consulted 8 October 2026.

\item \textbf{Unresolved original citation:} ``carbon-leakage literature.'' No author, title or year was supplied, so no unique source can be assigned.

\end{enumerate}

\subsection*{Common conventions: Source mathematical conventions}

Unless an entry states otherwise, agent and firm sets are finite. State and action spaces are measurable, strategies and outcome maps are measurable, probability laws are specified, and expectations exist for the targets under discussion. Infinite-horizon discount factors lie in $(0,1)$ and discounted objectives are finite. Norms, tolerances and complexity input sizes must be specified for computational comparisons. Constants or primitives that appear locally are inputs, not empirical facts asserted by this catalogue.

\textbf{Identification.} A statistical model $m\in\mathcal M$ implies an observable law $P_m$ and target $\theta(m)$. At law $P$, its identified set is
\[
\Theta(P;\mathcal M)=\{\theta(m):m\in\mathcal M,\ P_m=P\}.
\]
Point identification holds when this set is a singleton, or equivalently when $P_{m_1}=P_{m_2}$ implies $\theta(m_1)=\theta(m_2)$ for all compatible models. Sharp bounds mean bounds attained, or approached when the boundary is not attained, by compatible models. Writing this set is a definition, not a solution: the substantive task is to establish informative restrictions and derive or estimate the set. An unrestricted-confounding model may give only trivial bounds.

\textbf{Inference.} Identification, estimability and valid uncertainty quantification are separate questions. Fix $\alpha\in(0,1)$. A confidence procedure $C_n$ over a declared law class $\mathcal P$ has asymptotic uniform coverage at level $1-\alpha$ when
\[
\liminf_{n\to\infty}\inf_{P\in\mathcal P}\Pr_P\{\theta(P)\in C_n\}\geq1-\alpha.
\]
For set-identified targets the coverage event must be defined separately, for example coverage of the full identified set. If an entry uses identification-indexed models instead of uniquely indexed laws, the same coverage convention is applied over those models. Pointwise limits do not establish uniform validity, and asymptotic validity does not guarantee finite-sample precision. Sample sizes, dependence structures and weak-identification sequences are part of each application.

\textbf{Counterfactuals.} $Y_i(a)$ is a potential outcome under a specified intervention, including its timing, implementation and versions. It is not $Y_i$ conditional on voluntarily choosing $a$. A full intervention vector permits interference; shorthand scalar treatment notation does not silently impose no interference. Assignment, selection, attrition, anticipation and measurement must be specified. A claim about an instrument's local effect is not automatically a population policy effect.

\textbf{Economic equilibrium.} A counterfactual model includes best responses, constraints, feasibility, market clearing, state transitions and expectations consistent with the stated information. When several equilibria exist, either a declared selector is an input or the target is set-valued. Comparisons across policy regimes must use the stated selection convention. Reduced-form relationships are not presumed invariant after an equilibrium-changing intervention.

\textbf{Optimization and welfare.} Optimization is over the stated feasible set. If attainment is not established, $\sup$ or $\inf$ and $\varepsilon$-optimal choices replace an asserted maximizer or minimizer. For example, with value $V=\sup_{a\in\mathcal A}W(a)<\infty$, an $\varepsilon$-optimal set is $\{a\in\mathcal A:W(a)\geq V-\varepsilon\}$ for $\varepsilon>0$. Benefits, costs, risk and health or environmental outcomes can be aggregated only using declared units and conversion weights. Interpersonal utility weights, the treatment of fiscal transfers and foreign profits, discounting, and the behavioral welfare criterion are explicit inputs. A comparative-static derivative additionally requires existence and appropriate differentiability of the selected counterfactual.

\textbf{Sources.} Local labels [S1], [S2], and so forth restart in every question. Locators identify the relevant abstract, model, section or result at the level actually checked; a bibliographic or abstract-level verification is not represented as inspection of an entire proof. Working papers are distinguished from later publications and changing versions. Source summaries are paraphrased. All URL checks and editorial formulations refer to the audit date stated above.
% END ECONOMICS STATEMENT
\end{document}
